\documentclass[00-Trabajo_Final.tex]{subfiles}

\begin{document}
\renewcommand{\NombreDocumento}{Resumen}
\subfile{Plantilla/Portada_documento}
\renewcommand{\contentsname}{Índice \NombreDocumento}
\localtableofcontents
\newpage

A continuación se incluye el título y una breve descripción a modo de resumen en castellano, gallego e inglés, como requerido por la norma de formato de trabajo final.

\section{Título y resumen}

\textit{Escalabilidad y eficiencia en Inteligencia Artificial a través de algoritmos federados verdes}

\vspace{1em}
\noindent
En este Trabajo Final de Grado se ha abordado el desarrollo, análisis y experimentación con métodos de aprendizaje federado desde una perspectiva verde, en sintonía con la manifestación de interés del Programa Nacional de Algoritmos Verdes \cite{libro_PNAV} publicado por el Ministerio de Asuntos Económicos y Transformación Digital \cite{web_ministerio}.

Para ello se dio comienzo con un repaso por las principales ramas y estado del arte en Inteligencia Artificial y Aprendizaje Automático hasta llegar a las técnicas proporcionadas por una tendencia reciente conocida como Aprendizaje Federado.

Estas últimas se presentan como una alternativa que solventa las desventajas del entrenamiento centralizado corriente, que requiere grandes silos de información y un mayor consumo y tiempo para su entrenamiento. El aprendizaje federado es así distribuido, eficiente y privativo por naturaleza, por lo que es un incentivo para la convergencia de las nuevas tecnologías de la información como \textit{IoT} y las operativas propias de la \textit{Industria 4.0}. Al mismo tiempo son medio-ambientalmente sostenibles, ya que requieren consumos de energía menores que su contrapartida centralizada.

El propósito ha sido así el de experimentar con técnicas de aprendizaje federado, en especial contribuir y mejorar, bajo nuevos escenarios, el método de entrenamiento regularizado para redes neuronales de una sola capa que ha sido desarrollado por la UDC en entornos federados (\textit{FedROLANN}) \cite{paper_FEDROLANN} y la implementación de estos algoritmos en dispositivos embebidos y simulados.

Como tal, se ha procedido a aplicar métodos de ensamblaje a dicho algoritmo para verificar su rendimiento contra la versión original y, por último, se ha usado como pieza en un sistema formado junto con otro algoritmo federado de agrupamiento de características similares (\textit{k-FED}) \cite{paper_kFED}, con el objetivo de lograr, entre ambos, un algoritmo más potente y de mayor capacidad de representación que el original.\\\\En ambos casos se lograron con éxito dichas metas.

\section{Título e resumo}
\textit{Escalabilidade e eficiencia en Intelixencia Artificial a través de algoritmos federados verdes}

\vspace{1em}
\noindent
Neste Traballo Final de Grao abordouse o desenvolvemento, análise e experimentación con métodos de aprendizaxe federada dende unha perspectiva verde, en sintonía coa manifestación de interese do Programa Nacional de Algoritmos Verdes \cite{libro_PNAV} publicado polo Ministerio de Asuntos Económicos e Transformación Dixital \cite{web_ministerio}.

Para isto deuse comezo cun repaso polas principais ramas e estado da arte en Intelixencia Artificial e Aprendizaxe Automática ata chegar ás técnicas proporcionadas por unha tendencia recente coñecida como Aprendizaxe Federada.

Estas últimas preséntanse como unha alternativa ás desvantaxes do adestramento centralizado corrente, que require grandes silos de información e un maior consumo e tempo para o seu adestramento. A aprendizaxe federada é así distribuída, eficiente e privativa por natureza, polo que é un incentivo para a converxencia das novas tecnoloxías da información como \textit{IoT} e as operativas propias da \textit{Industria 4.0}. Ao mesmo tempo son medio-ambientalmente sostibles, xa que requiren consumos de enerxía menores que a súa contrapartida centralizada.

O propósito foi así o de experimentar cas técnicas de aprendizaxe federada, en especial contribuír e mellorar, baixo novos escenarios, o método de adestramento regularizado para redes neuronais dunha soa capa que foi desenvolto pola UDC en contornas federadas (\textit{FedROLANN}) \cite{paper_FEDROLANN} e a implementación destes algoritmos en dispositivos embebidos e simulados.

Como tal, procedeuse a aplicar métodos de ensamblaxe ao devandito algoritmo para verificar o seu rendemento contra a versión orixinal e, por último, usouse como peza nun sistema formado xunto outro algoritmo federado de agrupamento de características similares (\textit{k-FED}) \cite{paper_kFED}, co obxectivo de lograr, entre ambos, un algoritmo máis potente e de maior capacidade de representación que o orixinal.\\\\En ambos os casos lográronse con éxito ditas metas.

\section{Title and summary}
\textit{Scalability and efficiency in Artificial Intelligence through federated green algorithms}

\vspace{1em}
\noindent
This Final Degree Project has addressed the development, analysis and experimentation of federated learning methods within a green perspective, in line with the manifestation of interest of the National Green Algorithms Program \cite{libro_PNAV} published by the Ministry of Economic Affairs and Digital Transformation \cite{web_ministerio}.

To do so, we started with a review of the main fields and state-of-the-art methods in Artificial Intelligence and Machine Learning up to the techniques provided by a recent trend known as Federated Learning.

The latter is presented as an alternative that overcomes the disadvantages of the current centralized training methods, which require large silos of data and a greater energy-consumption and time for its training. Federated learning is thus distributed, efficient and privative by nature, making it an incentive for the convergence of new information technologies such as \textit{IoT} and \textit{Industry 4.0} operations. At the same time they are environmentally sustainable, as they require lower energy consumptions than their centralized counterpart.

The purpose has thus been to experiment with federated learning techniques, and in particular to contribute and improve, under new scenarios, the regularized training method for single-layer neural networks that has been developed by the UDC for federated environments (\textit{FedROLANN}) \cite{paper_FEDROLANN}, along with the implementation of these algorithms in embedded and simulated devices.

As such, ensemble methods have been applied to this algorithm to verify its performance against the original version and, finally, it has been used as a building block for a greater system formed together with another federated clustering algorithm of similar characteristics (\textit{k-FED}) \cite{paper_kFED}; with the aim of achieving, using both, a more powerful algorithm with a higher representation capacity than the original one. Both cases were successfully achieved.

\cleardoublepage
\end{document}